\ssr{ЗАКЛЮЧЕНИЕ}

В ходе выполнения курсовой работы была разработана программная система для моделирования световых паттернов статических декоративных гирлянд, обеспечивающая физически обоснованный расчет распределения света и интерактивную визуализацию результатов.

В процессе работы были решены следующие задачи:

\begin{enumerate}
    \item разработана математическая модель распространения света от точечных источников с учетом преломления в защитных колпачках (закон Снеллиуса, аппроксимация Френеля, поглощение и фильтрация, стохастическое рассеивание) и реализован алгоритм трассировки лучей для расчета распределения освещенности на диффузных поверхностях (модель Ламберта);
    
    \item создана система предварительного расчета освещения с формированием карт освещенности и поддержкой стохастического рассеивания для имитации матовых материалов колпачков;
    
    \item разработан графический пользовательский интерфейс с инструментами создания паттернов гирлянды (точка, линия, дуга, круг) и редактирования параметров источников света и защитных колпачков;
    
    \item реализовано интерактивное управление камерой и визуализация результатов расчета в реальном времени на основе предрассчитанных карт освещенности;
    
    \item проведен анализ производительности системы и исследование зависимости времени расчета от количества источников света, объектов сцены и разрешения карт освещенности.
\end{enumerate}

Разработанный инструмент позволяет в реальном времени визуализировать каустики, на расчет которых в универсальных рендерерах требуется в 3--5 раз больше времени на настройку материалов. Лимитирующим фактором производительности является этап генерации лучей (80\% времени), запись в текстурные атласы вносит менее 12\% задержки. Для интерактивной работы оптимальным является использование 5--10 источников света и 3--5 объектов, что обеспечивает время расчета менее 5 секунд.

Текущая версия не учитывает вторичные отражения от стен, что планируется реализовать через внедрение Light Propagation Volumes. Система может быть использована для предварительной визуализации декоративного освещения в архитектурном дизайне, планирования световых инсталляций и образовательных приложений для демонстрации принципов геометрической оптики.

\clearpage
